\chapter{Données et catalogue de référence}
\label{annexe:donnees}

\section{Catalogue des 100 preuves attendues}

Le catalogue de référence pour la mesure de couverture a été
construit par lecture du \emph{Guide d'homologation ANSSI} et de
l'\emph{ISO/IEC 27002:2022}. Chaque preuve est classée~:

\begin{itemize}
  \item \textbf{A}~: automatisable (collectable par un script
    depuis une API ou un fichier).
  \item \textbf{B}~: semi-automatisable (collectable mais avec
    intervention humaine pour valider la pertinence).
  \item \textbf{C}~: manuelle (formation, sensibilisation,
    revue d'analyse de risque, signature d'un document).
\end{itemize}

\section{Décomposition par étape ANSSI}

\begin{table}[h]
\centering
\caption{Répartition des 100 preuves par étape ANSSI.}
\label{tab:catalogue}
\small
\begin{tabular}{lrrrrr}
\toprule
\textbf{Étape ANSSI} & \textbf{Total} & \textbf{A} & \textbf{B} & \textbf{C} & \textbf{Couvert} \\
\midrule
1. Cadrage & 6 & 0 & 1 & 5 & 0 \\
2. Périmètre fonctionnel & 5 & 0 & 2 & 3 & 1 \\
3. Acteurs et rôles & 7 & 1 & 2 & 4 & 1 \\
4. Cartographie SI & 18 & 14 & 3 & 1 & 14 \\
5. Analyse de risque EBIOS RM & 8 & 0 & 3 & 5 & 1 \\
6. Mesures de sécurité & 28 & 19 & 6 & 3 & 16 \\
7. Audits et tests & 16 & 11 & 4 & 1 & 9 \\
8. Décision d'homologation & 0 & 0 & 0 & 0 & 0 \\
9. Maintien dans le temps & 12 & 7 & 3 & 2 & 5 \\
\midrule
\textbf{Total} & \textbf{100} & \textbf{52} & \textbf{24} & \textbf{24} & \textbf{47} \\
\bottomrule
\end{tabular}
\end{table}

Note~: l'étape~8 (décision d'homologation) est par nature un acte
signé par l'AQSSI, non automatisable. Elle ne compte pas dans le
score.

\section{Exemple de preuves par catégorie}

\subsection{Catégorie A (automatisable)}

\begin{itemize}
  \item Liste des règles NSG et leur état \texttt{Allow}/\texttt{Deny}.
  \item Liste des CVE \texttt{HIGH} et \texttt{CRITICAL} sur les
    images Docker en production.
  \item Liste des comptes AD inactifs (\texttt{lastLogon}~>
    90~jours).
  \item Liste des règles Wazuh activées par catégorie.
  \item État des agents Wazuh (nombre, dernière connexion).
\end{itemize}

\subsection{Catégorie B (semi-automatisable)}

\begin{itemize}
  \item Cartographie des flux entre applications (peut être
    extraite des logs nginx, mais demande validation manuelle).
  \item Inventaire des comptes à privilèges (extractible LDAP,
    mais la qualification \emph{«~privilégié~»} demande revue).
  \item Liste des journaux activés sur les composants critiques
    (à corréler avec le catalogue attendu).
\end{itemize}

\subsection{Catégorie C (manuelle)}

\begin{itemize}
  \item Compte-rendu des sessions de sensibilisation des
    utilisateurs.
  \item Procédure de gestion des incidents et son test annuel.
  \item Compte-rendu de l'atelier EBIOS RM atelier~3
    (scénarios stratégiques).
  \item Signature de l'AQSSI sur la décision d'homologation.
\end{itemize}

\section{Données de simulation pour l'audit drift}

La simulation rétroactive de 180 jours utilise un modèle de
changement calibré sur le journal git du lab~:

\begin{table}[h]
\centering
\caption{Taux de changement mensuel par catégorie (calibré sur
le lab).}
\label{tab:churn}
\small
\begin{tabular}{lr}
\toprule
\textbf{Catégorie} & \textbf{Changements / mois} \\
\midrule
Règles NSG & 2--4 \\
CVE (apparition + correction) & 8--15 \\
Comptes AD & 1--3 \\
Règles Wazuh & 1--2 \\
\bottomrule
\end{tabular}
\end{table}

Ces taux sont volontairement bas (lab modeste). Sur un ENT en
production, ils seraient probablement multipliés par 5 à 10, ce
qui rendrait l'écart entre dossier statique et dossier continu
encore plus marqué.
